% !TEX program = xelatex
\documentclass[10pt,UTF8]{article}
\usepackage{design_ASC}
\usepackage{enumerate}
\usepackage{amsmath}
\usepackage{mathrsfs}
\usepackage{commutative-diagrams}
\usepackage{thmtools}
\usetikzlibrary{arrows}
\usetikzlibrary{arrows.meta}
\usepackage{amsthm}
\theoremstyle{definition}
\usepackage{mathdots}
\usepackage[slantfont, boldfont]{xeCJK}

\renewcommand{\qed}{\rightline{$\qedsymbol$}}

\renewcommand{\emptyset}{\varnothing}
\renewcommand{\d}{{\rm d}}
\renewcommand{\i}{{\rm i}}
\newcommand{\id}{{\rm id}}
\renewcommand{\div}{{\rm div}}
\newcommand{\rot}{{\rm rot}}
\newcommand{\D}{{\rm D}}
\newcommand{\ad}{{\rm ad}}
\newcommand{\Aut}{{\rm Aut}}
\newcommand{\Int}{{\rm Int}}
\newcommand{\End}{{\rm End}}
\newcommand{\Ker}{{\rm Ker}}
\newcommand{\Gal}{{\rm Gal}}
\renewcommand{\Im}{{\rm Im}}
\newcommand{\Irr}{{\rm Irr}}
\newcommand{\Frac}{{\rm Frac}}
\renewcommand{\char}{{\rm char}}
\newcommand{\gen}[1]{\langle #1 \rangle}

\newcommand{\ent}{\text{ent}}
\newcommand{\iden}{\text{id}}
\renewcommand{\R}{\mathcal{R}}
\renewcommand{\N}{\text{N}}
\newcommand{\T}{\text{T}}
\newcommand{\gl}{\text{GL}}
\renewcommand{\o}{\text{O}}
\renewcommand{\so}{\text{SO}}
\newcommand{\mso}{\text{MSO}}
\renewcommand{\hom}{\text{Hom}}
\newcommand{\spm}{\text{Spm}}
\newcommand{\M}{\text{M}}
\newcommand{\Q}{\mathbb{Q}}
\renewcommand{\k}{\Bbbk}
\renewcommand{\Z}{\mathbb{Z}}
\renewcommand{\tr}{{\rm tr}}
\newcommand{\sgn}{{\rm sgn}}
\renewcommand{\mod}[1]{{\rm \ mod}\  #1}
\renewcommand{\exp}[1]{ {\rm exp}  \left(#1 \right)}
\newcommand{\e}{{\rm e}}
\newcommand{\ord}{{\rm ord}}
\newcommand{\jie}{~\\ \textbf{解}}
\newcommand{\genlegendre}[4]{
  \genfrac{(}{)}{}{#1}{#3}{#4}
  \if\relax\detokenize{#2}\relax\else_{\!#2}\fi
}
\newcommand{\legendre}[3][]{\genlegendre{}{#1}{#2}{#3}}
\newcommand{\dlegendre}[3][]{\genlegendre{0}{#1}{#2}{#3}}
\newcommand{\tlegendre}[3][]{\genlegendre{1}{#1}{#2}{#3}}

\newtheorem{definition}{Definition}[section]
\newtheorem{theorem}[definition]{Theorem}
\newtheorem{lemma}[definition]{Lemma}
\newtheorem{proposition}[definition]{Proposition}
\newtheorem*{example}{Example}
\newtheorem*{remark}{Remark}
\let\proof\relax
\newtheorem*{proof}{Proof}


\setlength\parindent{2em} 
\title{The Decomposition of Non-Square Matrices} 
\author{}
\date{}



\begin{document}
\maketitle

\begin{abstract}
  The Bruhat decomposition of the reductive group is an important theorem in algebraic group theory. Givn that 
  a reductive group $G$, then it can be decomposed as the disjoint union of $BwB$ where $B$ is the Borel
  subgroup and $w$ is a element of the Weyl group. Hence, if we consider the quotient group $G/B$ with the torus action $T$ 
  on it, then it induces the cell decomposition of $G$ i.e. $G/B$ can be decomposed by the disjoint union of affine 
  subvariety. The cell decomposition is actually the Białynicki-Birula decomposition.

  In this paper, we firstly introduce the basic concepts of algebraic geometry and group theory. 
  Accordingly, we present the Bruhat decomposition of non-square matrices $M$ over an algebraically closed field $k$ and 
  verify the rationality of applying Białynicki-Birula decomposition on $M/B_m$. The main contribution of this paper is 
  \begin{enumerate}
    \item prove the Bruhat decomposition of non-square matrices;
    \item constructing the quasi-flag from $M/B_m$;
    \item prove that quasi-flag is a projective variety;
    \item prove the smoothness of quasi-flag;
    \item define a locally affine torus action on $M/B_m$.
  \end{enumerate}
  \textbf{keywords}: Non-square Matrices, Bruhat Decomposition, Błynicki-Birula Decomposition, Torus Action.
\end{abstract}


\begin{abstract}
    约化群的 Bruhat 分解在代数群理论中是十分重要的。对于任意的约化群 $G$，可以将其分解为若干个形如 $BwB$ 的不交并，其中
    $B$ 为约化群 $G$ 的博雷尔子群，$w$ 为外尔群中的元素。进一步考虑商群 $G/B$ 与商群上的代数环面作用，则有商群的胞腔分解，即
    可以分解为若干个仿射子簇。该胞腔分解实际上是 $G/B$ 的 Białynicki-Birula 分解。

    在本为中，我们首先介绍了代数几何与代数群的基本概念。基于此，我们给出了非方矩阵 $M$ 在代数闭域上的 Bruhat 分解。
    我们还验证了 Białynicki-Birula 分解在 $M/B_m$ 的合理性。本文的主要贡献为
    \begin{enumerate}
      \item 证明了非方矩阵的 Bruhat 分解；
      \item 定义了一个与 $M/B_m$ 等价的数学概念：仿旗 (Quasi-flag)；
      \item 证明了仿旗是一个投影簇；
      \item 证明了仿旗的光滑性；
      \item 在仿旗上定义了一个局部仿射性的代数环面作用。
    \end{enumerate}
    \textbf{关键词}: 非方矩阵, Bruhat 分解, Błynicki-Birula 分解, 代数环面作用.
\end{abstract}
\newpage
\tableofcontents
\newpage

\setcounter{page}{1}
\section{Introduction}

In this paper we study two relative and classical decomposition: the Bruhat decomposition
and Białynicki-Birula decomposition and try to apply it to non-square matrices.


Recall that the Bruhat decomposition is defined on algebraic group which has good properties such as smoothness and 
so on. Those good properties give a ample structure of algebraic group. For example, for any algebraic group, the lie algebra can 
be naturally defined, which leads us define many similar concepts in algebraic group theory. 
It contributes to the proof of Bruhat decomposition and understanding it's geometry meaning.
However, it is really hard for us to do same thing in non-square matrices. 

There are many works try to entend the Bruhat decomposition. For example, Pushkar\cite{pushkar2020enhanced} 
defines the enhanced Bruhat decomposition on non-square matrices by fixed the flag of a linear space, which is 
called enhanced linear space.

Another interesting theory is developed Putcha and Renner \cite{putcha1988linear,renner2005linear}, the algebraic semigroup theory.
They study the structure of algebraic moniod i.e. algebraic variety with the structure of moniod and try to extend the arguments of the 
algebraic group theory. One of the most important works is that the anologe of Bruhat decomposition on the reductive moniod, which was
presented by Renner\cite{renner1986analogue}. Renner proved that for any reductive moniod $M$, there are a decomposition
\begin{align*}
  M = \coprod_{x\in\mathcal{R}} BxB
\end{align*}
where $B$ denotes the Borel subgroup of unit group and the $\mathcal{R}$ is the Renner moniod of $M$. 
The Renner moniod is defined by a Zariski closed of Weyl group of unit group. In the sense of square matrices,
the Renner moniod is the rook matrices which has at most one non-zero element in each rows and columns and the non-zero 
element must be 1. It is a fantastic method which use the Zariski closed to extend the concepts from algebraic group to
moniod. 

For the decomposition of a non-square matrices, there are basic views to consider it and relative concepts
\begin{enumerate}
  \item as a algebraic variety and consider the decomposition of variety
  \item as a algebraic variety with a group action, and extend the decomposition of group to itself
  \item as a algebraic variety with a algebra structure, do similar thing with algebraic moniod.
\end{enumerate}

Our discussion of this paper is basically based on the first idea which is the simplest. 
For second idea, we view the non-square matrices as a variety with a algebraic group action. And this idea
is not really clear now. May be we can refer to \cite{jelisiejew2019bialynicki}, which define a moniod action as 
a extended group action.
For the last idea, may be we can embed the non-square matrices into a algebraic group or semigroup and consider the decomposition of it.

This paper is divided by three parts. On the first part, we give some basic concepts and important theorem of algebraic geometry and algebraic group theory.
On the second part, we prove the Bruhat decomposition of non-square matrices.
At last, we want to apply Białynicki-Birula decomposition to non-square matrices quoient by a Borel subgroup and verify its rationality.

This paper is based on the algebraic geometry and algebraic group theory, which is extremely hard for me to understand.
Therefore, the proof of some theorem may have some errors in details or not relatively strict, which 
needs some refinement and correction in my future study.




\newpage

\section{Preliminaries}

In this chapter, we give a brief introduction of some basic concepts to algebraic group theory.
For details, we refer to \cite{roman2005advanced,axler1997linear} for linear algebra,\cite{hartshorne2013algebraic,bosch2022algebraic,harris2013algebraic} for algebraic geometry,
\cite{milne2017algebraic,springer1994linear,geck2013introduction} for algebraic group theory.



.


\subsection{Some notations}
In this section, we introduce some notations used in this paper.

For any finite set $A$, $|A|$ denotes the order of set i.e. the number of elements in $A$.

$\mathbb{N}_{\le n}$ denote the integer set $\{0,\cdots,n\}$, $\mathcal{N}_{d,n} := \{A\subseteq \mathbb{N}_{\le n}||A| = d\}$.
Similarly, $\mathbb{N}^*_{\le n} := \{1,\cdots,n\}$ and $\mathcal{N}^*_{d,n} := \{|A| = d|A\subseteq \mathbb{N}^*_{\le n}\}$.

For matrix $A$, $\text{ent}_{ij}(A)$ means $ij$ entry of matrix, $\text{ent}_{* j}(A)$ denotes the $j$th column of $A$ and 
$\text{ent}_{i*}(A)$ denotes the $i$th row of $A$. Suppose $A$ has dimension $n\times m$, and $I\subseteq \mathbb{N}^*_{\le n}$ and $J\subseteq \mathbb{N}^*_{\le m}$, $\Delta_{I,J}(A)$ denotes the $IJ$ minors of $A$.
 
\begin{definition}
  The lexicographic order $L_d$ of $\mathcal{N}^*_{d,n}$ is given by
\begin{align*}
  L(I) = \sum_{i=1}^d \sum_{j=1}^{I_i}\binom{n-j}{d-i}.
\end{align*}
where $I=\{I_1<\cdots<I_d\}$.
\end{definition}



\subsection{Linear algebra and basic geometry}

Let $V$ be the $n$-dimensional vector space over field $k$ with a basis $\{v_1,\cdots,v_n\}$.






\subsubsection{Affine and projective space}
\begin{definition}
  The $n$-dimensional \textbf{affine space} over $k$ is the affine structure $\left\langle k^n,k^n,+ \right\rangle$, where $+:k^n\times k^n\to k^n$ is a transitive and free action, denoted by $\mathbb{A}^n_k$.
\end{definition}



\begin{definition}
  The \textbf{projective space} of vector space $V$ is defined by 
  \begin{align*}
    \mathbb{P}(V) := V^*/k^*
  \end{align*}
  where $V^*=V\backslash \{0\},k^*=k\backslash \{0\}$. The $n$-dimensional projective sapce over $k$ is 
  \begin{align*}
    \mathbb{P}^n_k := \mathbb{A}_k^{n+1}\backslash\{0\}/k^*.
  \end{align*}
\end{definition}
By definition, we have $\mathbb{P}(V)=\mathbb{P}^{n-1}_k$.
\begin{remark}
  We briefly write $\mathbb{P}^n$ and $\mathbb{A}^n$ as $\mathbb{P}_k^n$ and $\mathbb{A}_k^n$, if we fix the base field $k$.
\end{remark}

\begin{definition}
  For point $p\in \mathbb{P}^n$, we give the \textbf{ homogeneous coordinates} as 
  \begin{align*}
    p = [p_0:p_1:\cdots:p_n] 
  \end{align*}
  where $(p_0,\cdots,p_n)\in \pi^{-1}(p)$ and $\pi:\mathbb{A}^{n+1}\to \mathbb{P}^{n}$ is natural homomorphism.
\end{definition}

For any $p\in \mathbb{P}^n$, the homogeneous coordinate is not uniquely determined. For example,
\begin{align*}
  [p_0:p_1:\cdots:p_n]  = [\lambda p_0:\lambda p_1:\cdots:\lambda p_n]
\end{align*}
However, if we fix one non-zero element in homogeneous coordinates then the other element is determined. 
This is the idea of affine chart and affine coordinate of projective space.
\begin{definition}
  The standard \textbf{affine chart} of projective space $\mathbb{P}^n$ is the set $\{U_i\}_{i=0}^n$ where
  \begin{align*}
    U_i := \{[p_0:\cdots:p_n]|p_i\neq 0\} 
  \end{align*}
  with the coordinate function $\phi_i:U_i \to \mathbb{A}^{n}$ such that
  \begin{align*}
    [p_0:\cdots:p_n] \mapsto \left(\frac{p_0}{p_i},\cdots,\frac{p_{i-1}}{p_i},\frac{p_{i+1}}{p_i},\cdots,\frac{p_{n}}{p_i}\right)
  \end{align*}
  For $p\in U_i$, we say $\phi_i(p)$ is the \textbf{affine coordinate} of $p$ in the affine chart $U_i$.
\end{definition}

The coordinate function $\phi_i$ give the set isomorphism of $U_i\cong \mathbb{A}^n$. Hence, $\mathbb{P}^n$ can be viewed as 
a union of $n+1$ copies of $\mathbb{A}^n$.
 
\subsubsection{Exterior algebra of vector space}
\begin{definition}
  The \textbf{$l$th tensor power} of $V$ is given by the $l$ times tensor product of $V$:
  \begin{align*}
    T^l(V):=\underbrace{V\otimes \cdots \otimes V}_l.
  \end{align*}
  Specifically, $T^0(V)=k$.
  The \textbf{tensor algebra} of $V$ is the direct sum of all tensor power of $V$:
  \begin{align*}
    T(V) := \bigoplus_{l=0}^{\infty} T^l(V).
  \end{align*}
\end{definition}





\begin{definition}
  The \textbf{exterior product}  of vector space $V$ is the construction on any finite order set of vector:
  \begin{align*}
    \{w_1,\cdots,w_l\} \stackrel{\text{construction}}{\longrightarrow} w_1\wedge \cdots \wedge w_l
  \end{align*}
  such that for any $w\in V$, $w \wedge w = 0$ and satisfy associative 
  \begin{align*}
    (w_1\wedge w_2) \wedge w_3 = w_1\wedge (w_2 \wedge w_3)
  \end{align*}
  and bilinear condition 
  \begin{align*}
    &(w_1+w_2) \wedge w_3 = w_1\wedge w_3 +w_2\wedge w_3\\
    &w_3 \wedge (w_1+w_2)  = w_3\wedge w_1 +w_3\wedge w_2.
  \end{align*}
\end{definition}

\begin{definition}
  The \textbf{exterior algebra} of $V$ is the quotient algebra of the tensor algebra $T(V)$:
  \begin{align*}
    \bigwedge(V) := T(V)/(\{x\otimes x|x\in V\})
  \end{align*}
  with the exterior product.
  As a vector space, it can be direct sum of all exterior power of $V$:
  \begin{align*}
    \bigwedge(V) = \bigoplus_{i=0}^{\infty} \bigwedge^i(V).
  \end{align*}

  The \textbf{$l$th exterior power} of $V$ is the linear subspace of exterior algebra of V such that
  \begin{align*}
    \bigwedge^l(V) := \text{span}(v_{i_1}\wedge\cdots\wedge v_{i_l}|\text{$\{i_j\}_{j=1}^l\in \mathcal{N}^*_{l,n}$ and $i_1\le\cdots\le i_l$}).
  \end{align*}
  Specifically, $\bigwedge^0(V)=k$.
  The \textbf{$l$-sum exterior power} is the linear subspace of exterior algebra where 
  \begin{align*}
    \underline{\bigwedge^l}(V): = \bigoplus_{i=0}^{l} \bigwedge^i(V)
  \end{align*}
\end{definition}

\begin{definition}
  For any $w\in \bigwedge^l(V)$, by definition, we have 
\begin{align*}
  w = \sum_{I\in \mathcal{N}^*_{l,n}} a_Iv_I
\end{align*}
where $v_I =v_{i_1}\wedge \cdots \wedge v_{i_l}$ if $I=\{i_1,\cdots,i_l\}$ and $a_I\in k$ called Plücker coordinate. Specially, when $l=0$, we let the Plücker coordinate of $w\in k$ be its own, denoted by
$a_{\emptyset} $.
\end{definition}





Next,we claim some properties of exterior product.
\begin{proposition}
  Consider the exterior product over $V$, then \label{252}
  \begin{enumerate}[a)]
    \item $w\wedge v + v\wedge w = 0$ for all $w,v\in V$
    \item $w_1\wedge\cdots\wedge w_d =0$ if and only if $\{w_1,\cdots,w_d\}$ is not linearly independent.
    \item If $\{w_1,\cdots,w_n\}$ and $\{u_1,\cdots,u_n\}$ is linearly equivalence then 
    \begin{align*}
      w_1\wedge\cdots \wedge w_n = \det(B) u_1\wedge\cdots \wedge u_n 
    \end{align*}
    where $B$ is the base transformation matrix. 
  \end{enumerate}
\end{proposition}


\begin{definition}
  Suppose $w\in \bigwedge^l(V)$ and $v\in V$. We say $v$ \textbf{divides} $w$ if there exist $w'\in\bigwedge^{l-1}(V)$ such that $w=w'\wedge v$.
\end{definition}
\begin{proposition}
  Let $V,v$ and $w$ be as above. $v$ divides $w$ if and only if $v\wedge w=0$
\end{proposition}
\begin{definition}
  We say $w\in \bigwedge^l(V)$ is totally decomposable if there exists $\{w_1,\cdots,w_l\}\subseteq V$ such that 
  \begin{align*}
    w = w_1\wedge \cdots \wedge w_l.
  \end{align*}
\end{definition}
Not all elements in exterior power is totally decomposable. We give the counter example as below.
\begin{example}
  Consider $e_1\wedge e_2 + e_2\wedge e_3\in \bigwedge^2(\mathbb{R}^3)$ is not totally decomposable.
\end{example}
\begin{proposition}
  For any $w\in\bigwedge^l(V)$
  The set $\{v\in V|w\wedge v=0\}$ is a linear subspace of $V$.
\end{proposition}

\begin{definition}
  Define a map for $l$th exterior power
  \begin{align*}
    \varphi_l:  \bigwedge^{l}(V) &\to \hom\left(V,\bigwedge^{l+1}(V)\right) \\
    \omega &\mapsto \varphi_l(\omega) 
  \end{align*}
  where $\varphi_l(w)(v) = w\wedge v$ for $v\in V$. 
\end{definition}

As usual, we choose the standard base $\{e_I\}_{I\in \mathcal{N}^*_{l,n}}$ of $\bigwedge^{l}(V)$ associated with the lexicographic order $L_l$.
Suppose $w\in \bigwedge^{l}(V)$ have the Plücker coordinate $\{a_I\}_{I\in \mathcal{N}^*_{l,n}}$.
Let $D_l(w)$ be the $\tbinom{n}{l+1}\times n$ matrices represent the homomorphism $\varphi_l(w)$.
the $ij$-entry is given by
\begin{align*}
  d_{ij} = 
  \begin{cases}
    0 & j\notin L_{d+1}^{-1}(i)\\
    (-1)^{\tau_{ij}}a_I & \text{otherwise}
  \end{cases}
\end{align*}
where $L_d$ is the lexicographic order of $\{1,\cdots,d\}$, $I =L_{d+1}^{-1}(i)/\{j\} $ and 
$\tau_{ij}= |\{I_k>j|I_k\in I\}|$.

We have already given a several concepts on $l$th exterior power. By the construction of $l$sum exterior power, we can extend such concepts to $l$sum exterior power.
\begin{definition}
  For given $l$sum exterior power $\underline{\bigwedge}^l(V)$ of $V$
  \begin{enumerate}[a)]
    \item Define the projection $\pi_i:\underline{\bigwedge}^l(V) \to \bigwedge^i(V)$ for $i\le l$ as a linear transformation such that 
    \begin{align*}
      w = \sum_{|I|\le l,I\subseteq\{1,\cdots,n\}} a_Iv_I \mapsto \sum_{|I|= i,I\subseteq\{1,\cdots,n\}} a_Iv_I 
    \end{align*}
    \item For $w\in\underline{\bigwedge}^l(V)$, define the map $\tilde{\varphi}_i:\underline{\bigwedge}^l(V) \to \hom\left(V,\bigwedge^{i+1}(V)\right)$ such that $\tilde{\varphi}_i = \varphi_i \circ \pi_i$.
    \item Define $\tilde{D}_i(w) =D_i\circ\pi_i(w) $.
  \end{enumerate}
\end{definition}


\begin{remark}
  Sometimes we remove the tilde of notation defined on $l$sum exterior power for convinence.
\end{remark}

\begin{proposition}
  Let $\{w_i\}_{i=1}^l$ is a set of vector in $V$. Suppose
  \begin{align*}
    w_1\wedge \cdots \wedge w_l = \sum_{I\in \mathcal{N}_{l,n}} a_Iv_I
  \end{align*}
  where $a_I\in k$, $v_I = v_{I_1}\wedge\cdots \wedge v_{I_l}$ and $I=\{I_1<\cdots<I_l\}$. Let $W=[w_1,\cdots,w_l]$ as a matrix, then
  \begin{align*}
    a_{I} = \det(\Delta_{I,I}).
  \end{align*}
  where $\Delta$ denote the minors.
\end{proposition}



\subsection{Some algebraic geometry and algebraic group}

Let $k$ be algebraically closed field and consider the polynomial ring $k[X_1,X_2,\cdots,X_n]$, where the $X_i$ are independent variables. 



\begin{definition}
  $X\subseteq k^n$ is \textbf{an affine variety} if $X$ is the set of common zeros of a finite collection of polynomials $S\subseteq k[X_1,\cdots,X_n]$.
\end{definition}

\begin{definition}
  The topology on $k^n$ induced by which affine varieties is closed is called \textbf{Zariski topology}.
\end{definition}


\begin{definition}
  \textbf{Affine space} in algebraic geometry is defined by 
  \begin{align*}
    \mathbb{A}^n := \text{spm}(k[X_1,\cdots,X_n]) 
  \end{align*}
  where $\text{spm}$ means the set of maximal ideals endowed with Zariski topology.
\end{definition}


\begin{definition}
  An affine variety $X$ is \textbf{irreducible} if $X$ is irreducible with the induced subtopology i.e. it is not the union of two non-trivial closed subsets.
\end{definition}

\begin{definition}
  Let $X$ be the affine variety of $k^n$ and $I(X)=\{f\in k[X_1,\cdots,X_n]|\forall x\in X,f(x)=0\}$. Obviously, $I(X)$ is an ideal of $k[X_1,\cdots,X_n]$. 
  \begin{align*}
    k[X_1,\cdots,X_n]/ I(X)
  \end{align*}
  is called the \textbf{coordinate ring} or affine algebra of $X$. It can be denoted by $\mathcal{O}(X)$.
\end{definition}

\begin{definition}
  Let $X$ be an affine variety with coordinate ring $\mathcal{O}(X)$. A function $f:X\to k$ is \textbf{regular at point} $p\in X$ if there is an open neighbourhood $U$ and $g,h \in\mathcal{O}(X)$ such that $f=g/h$ on $U$.
  We say $f$ is regular if it is regular at every point of $X$.
\end{definition}


\begin{definition}
  Let $X\subseteq k^n$ and $Y\subseteq k^m$ be affine varieties. A map $\psi:X\to Y$ is a morphism if
  \begin{align*}
    \psi(x_1\cdots,x_n) = (\psi_1(x),\cdots,\psi_m(x))
  \end{align*}
  where $\psi_i \in \mathcal{O}(X)$ for $i=1,\cdots,n$ and $x=(x_1,\cdots,x_n)$. We also call this type of map as \textbf{regular map}.
  A regular map whose inverse is also regular is called \textbf{biregular}. The \textbf{isomorphism} between two affine variety is regular and biregular.
\end{definition}


\begin{definition}
  The Zariski topology on $\mathbb{P}^n$ is the topology whose open sets are the complements of 
  projective algebraic sets.
\end{definition}

\begin{definition}
  A projective variety is an irreducible algebraicset in $\mathbb{P}^n$ endowed with the Zariski topology. 
\end{definition}



\begin{definition}
  Let $G$ be an affine algebraic variety over $k$ with group structure. If multiplicative function $m:G\times G\to G$ and inverse function$\tau:G\to G $ where $\tau(x)=x^{-1}$ are regular maps, then we call
  $G$ \textbf{linear algebraic group}.
\end{definition}
\begin{definition}
  The multiplicative group is $\mathbb{G}_m=(k^*,\cdot)$.
\end{definition}
\begin{definition}
  \textbf{A morphism of algebraic groups} is a group homomorphsim and a morphism of varieties(regular map).
\end{definition}

\begin{definition}
  For a finite-dimensional $k$-vector space $V$, let $\gl_V$ denote linear automorphism group of $V$ and $G$ be is a algebraic group.
  A \textbf{representation} of $G$ is a homomorphsim $r:G\to \gl_V$.
\end{definition}

\begin{definition}
  A homomorphsim of algebraic groups $\chi:G\to \mathbb{G}_m$ is called a (rational) \textbf{character}. Conversely, homomorphsim $\lambda:\mathbb{G}_m\to G$ called a cocharacter of $G$.
\end{definition}

It is easy to find that a character $\chi$ of $G$ can induce a representation $r$ of $G$ by
\begin{align*}
  r(g)v = \chi(g)v \quad  \forall g\in G, v\in V.
\end{align*}

We also say a cocharacter as a one-parameter subgroup of $G$.
\begin{definition}
  An algebraic group $G$ over $k$ is a \textbf{torus} if it is isomorphic to a product of copies of $\mathbb{G}_m$.
\end{definition}


\begin{definition}
  Let $R(G)$ denote the radical of $G$ i.e. the largest connected sovable normal subgroup variety. We say $G$ is \textbf{semisimple} if $R(G) = e$.
\end{definition}



\begin{definition}
  Let $G$ be a connected group over $k$. unipotent radical $R_u(G)$ of $G$ is the largest connected normal unipotent subgroup of $G$. 
  $G$ is said to be \textbf{reductive} if $R_u(G)=e$.
\end{definition}

\begin{definition}
  Let $X$ be a separated algebraic scheme over $k$ and $\varphi:\mathbb{A}^1\backslash\{0\}\to X$ a regular map.
  If $\varphi$ extends to a regular map $\tilde{\varphi}:\mathbb{A}^1\to X$, then define
  \begin{align*}
    \lim_{t\to 0} \varphi(t)  = \tilde{\varphi}(0)
  \end{align*}
\end{definition}

\begin{definition}
  Let $G$ be an algebraic group and $\lambda:\mathbb{G}_m \to G$ a cocharacter of $G$. Then $\lambda$ defines an action of $\mathbb{G}_m$ on $G$ by
  \begin{align*}
    (t,g) \mapsto \lambda(t)g\lambda(t)^{-1}
  \end{align*}
  where $t\in \mathbb{G}_m,g\in G$, denoted by $t.g$.
\end{definition}
\begin{definition}
  Let $G$ be a algebraic group over $k$ and $\lambda$ a cocharacter of $G$, define
  \begin{align*}
    &P_G(\lambda) = \{g\in G|\lim_{t\to 0} t.g\text{ exists}\}\\
    &U_G(\lambda) = \{g\in P_G(\lambda)|\lim_{t\to 0} t.g = e\}\\
    &Z_G(\lambda) = P_G(\lambda)\cap P_G(-\lambda)
  \end{align*}
\end{definition}




\subsection{The Białynicki-Birula decomposition}

\begin{definition}
  Let $\mu:G\times X\to X$ be an action of a algebraic group on algebraic scheme $X$ over $k$.
  Define the \textbf{fixed subscheme} of $X$ as a largest closed subscheme $X^G$ of $X$ on which
  $G$ acts trivially.
\end{definition}



\begin{definition}
  An action of a torus $T$ on a scheme $X$ over $k$ is \textbf{locally affine} if $X$ 
  admits a covering by $T$-invarient open affine subschemes.
\end{definition}



\begin{theorem}
  Let $X$ be a smooth algebraic variety over $k$ equipped with a locally affine action of $\mathbb{G}_m$.
  Let $Z$ be a connected component of $X^{\mathbb{G}_m}$. There exist a unique smooth
    subvariety $X(Z)$ of $X$ such that
    \begin{align*}
      X(Z) = \{y\in X|\lim_{t\to 0}ty\text{ exists and lie in }Z\}
    \end{align*}
    and a unique regular map $\gamma_Z:X(Z) \to Z$ sending $y\in X(Z)$ to the $\lim_{t\to 0}ty\in Z$.
\end{theorem}

\begin{theorem}
  Let $X$ be a smooth connected scheme over $k$ with a locally affine action of $\mathbb{G}_m$ such that
  $X^{\mathbb{G}_m}$ is finite and constant. For each $x\in X^{\mathbb{G}_m}$, there is a unique
  smooth subscheme $X(x)$ of $X$ such that 
  \begin{align*}
    X(x) = \{y\in X|\lim_{t\to 0}ty\text{ exists and equals x}\}
  \end{align*}
\end{theorem}



\subsection{The Bruhat decomposition}



\begin{definition}
  Let $G$ be a reductive algebraic group over $k$ with maximal torus $T$. The Weyl group of $G$ is defined by 
  \begin{align*}
    W(G,T) = {N_G(T)}/{C_G(T)} = {N_G(T)}/{T}
  \end{align*}
  which is a finite algebraic group. 
\end{definition}

\begin{definition}
  A \textbf{Borel subgroup} $B$ of $G$ is a maximal connected solvable subgroup of $G$.
\end{definition}

\begin{theorem}
  Let $(G,T)$ be a reductive group over $k$ with maximal torus and $B$ a Borel subgroup of $G$ containing $T$. 
  There is Bruhat decomposition of $G$ that 
  \begin{align*}
     G = \coprod_{w\in W} BwB
  \end{align*} 
  where $W$ is Weyl group of $(G,T)$.
\end{theorem}
\begin{theorem}
  Using previous notation, for $w\in W$, let $e_w = wB/B$ be a fixed point in $G/B$ under the action of torus $T$, i.e. $e_w \in (G/B)^T$, there is a unique smooth subvariety $Y(w)$ of $G/B$ such that 
  \begin{align*}
    Y(w) = \{x\in G/B|\lim_{t\to 0}\lambda(t)\cdot x = e_w\}
  \end{align*}
  which is an affine space and we have a decomposition
  \begin{align*}
    G/B = \coprod_{w\in W} Y(w).
  \end{align*}
\end{theorem}



\subsection{Grassmannian, Plücker embedding and flag variety}
\begin{definition}
  Let ${G}_{d,V}$ be the set of $d$-dimensional $k$-vector subspace of $V$, named as Grassmannian.
\end{definition}


\begin{definition}
  Let $\mathcal{P}_{d,V}:G_{d,V} \to \mathbb{P}\left( \bigwedge^d(V)\right)$ denotes the Plücker embedding of $G_{d,V}$ such that
  \begin{align*}
    W \mapsto \overline{w_1\wedge\cdots \wedge w_d}
  \end{align*}
  where $\{w_1,\cdots,w_d\}$ is a basis of linear subspace $W$. Specifically, when $d=0$, the Plücker embedding $\mathcal{P}_0$ give the 
  \begin{align*}
    0 \mapsto \overline{k}
  \end{align*} 
  where $\overline{k}$ is the single point of $\mathbb{P}^0$.
\end{definition}


\begin{remark}
  Since we fix the vector space $V$ in following discussion, let $\mathcal{P}_d$ simply denotes $\mathcal{P}_{d,V}$.
\end{remark}

The Plücker embedding give a projective variety structure of Grassmannian. Therefore,
\begin{proposition}
  ${G}_{d,V}$ is a projective variety.
\end{proposition}


Here are a few propositions related to Plücker embedding.
\begin{proposition}
  $\mathcal{P}_d$ is injective. \label{2}
\end{proposition}

\begin{proposition}
  $\mathcal{P}_d$ is a closed embedding.
\end{proposition}

\begin{definition}
  The flag variety of linear space $V$ is 
  \begin{align*}
    \mathcal{F}(n_1,\cdots,n_r; V) = \{0\subseteq V_1 \subseteq \cdots \subseteq V_r \subseteq V|\dim(V_i) = n_i\}
  \end{align*}
  where $1\le n_1<\cdots<n_r\le n$ a sequence of integer.
\end{definition}
\begin{proposition}
  Flag variety is a smooth projective variety. \label{fv}
\end{proposition}


\subsection{Segre embedding}
\begin{definition}
  The Segre embedding $ \mathscr{S}: \mathbb{P}^m \times \mathbb{P}^n \to \mathbb{P}^{(n+1)(m+1)-1}$ is defined as the map
  \begin{align*}
    ([x_0:\cdots:x_n],[y_0:\cdots:y_m]) \mapsto [x_0y_0:x_0y_1:\cdots:x_iy_j:\cdots:x_ny_m]
  \end{align*}
  where $x_iy_j$ are taken in lexicographical order. Hence, we let $\mathscr{S}_2= \mathscr{S}$ and for $d>2$
  \begin{align*}
    \mathscr{S}_d:\mathbb{P}^{m_1}\times\cdots\times \mathbb{P}^{m_d} &\to \mathbb{P}^{(m_1+1)\cdots(m_d+1)-1}\\
    (p_1, \cdots ,p_d) &\mapsto \mathscr{S}(\mathscr{S}_{d-1}(p_1,\cdots,p_{d-1}),p_d)
  \end{align*}
\end{definition}

\begin{proposition}
  Segre embedding is a projective variety morphism.
\end{proposition}


\newpage

\section{The Bruhat decomposition of non-square matrices}


Recall that the Burhat decomposition of $\gl_n(k)$ is 
\begin{align*}
  \gl_n(k) = \coprod_{w\in W_n} B_nwB_n
\end{align*}
Where $B_n$ denotes the $n$-dimensional invertible upper triangular matrices as a Borel subgroup of $\gl_n(k)$ and $W$ denotes the $n$-dimensional permutation matrices
as a Weyl group. 

In this part, we claim that there are a similar decomposition of non-square matrices over $k$.
Let $M$ be the $n\times m$ matries over $k$.

\begin{definition}
  The $0-1$ matrix of $M$ is the $n\times m$ matries which entries are only 1 or 0.
\end{definition}

\begin{definition}
  Define the $n\times m$ rook matrices as 
  \begin{align*}
    \mathcal{R}_{n\times m} = \{g\in M| \text{$g$ is a 0-1 matrix at most one element in each row and column is 1}\}
  \end{align*} 
\end{definition}




Before giving the proof of Bruhat decomposition, we define some new concepts.
\begin{definition}
  For $n\times m$ matrix $A$, the $ij$ partial rank of matrix $A$ is defined by
  \begin{align*}
    r_{ij}(A) := \rank(\Delta_{I,J}(A))
  \end{align*}
  where $\Delta_{I,J}$ means $I,J$ minors and $I=\{i,\cdots,n\},J=\{1,\cdots,j\}$.
\end{definition}

\begin{lemma}
  Partial rank is invarient under $B_n$ and $B_m$ i.e. $r_{ij}(b_1A) = r_{ij}(A) = r_{ij}(Ab_2)$ for any $b_1\in B_n,b_2\in B_m$ and $A\in M$.  \label{pr}
\end{lemma}
\begin{proof}
  By the definition of matrix product, $b_1A$ means that we add to each row of A a linear combination of the rows further below. 
  This kind of row transformation is obviously not change the partial rank. It is similar for the column cases.

  \qed
\end{proof}

\begin{lemma}
  For any $w\in \mathcal{R}$, we have\label{rpr}
  \begin{align*}
    \ent_{ij}(w) = r_{ij}(w) -  r_{i(j-1)}(w) -  r_{(i+1)j}(w) + r_{(i+1)(j-1)}(w) \tag{$*$} 
  \end{align*}
\end{lemma}
\begin{proof}
  Suppose that $r_{ij}(w) -  r_{i(j-1)}(w) =0$ then $\ent_{lj}(w)=0$ for any $i\le l\le n$ and 
  $r_{(i+1)j}(w) - r_{(i+1)(j-1)}(w)=0$.
  If $r_{ij}(w) -  r_{i(j-1)}(w) = 1$ and $\ent_{ij}(w)=0$, then exist $i<l\le n$ such that $\ent_{lj}(w)\neq 0$. 
  As a consequence, $r_{(i+1)j}(w) - r_{(i+1)(j-1)}(w)=1$.
  If $r_{ij}(w) -  r_{i(j-1)}(w) = 1$ and $\ent_{ij}(w)=1$, then for all $i<l\le n$,  $\ent_{lj}(w)= 0$.
  Therefore, $(*)$ holds.

  \qed
\end{proof}


\begin{theorem}
  For $n,m\in \mathbb{N}^*$, $\mathcal{R}$ denotes the $n\times m$ rook matrix let $B_n$ denotes the invertible triangular matrices over $k$. Then
  we have decomposition 
  \begin{align*}
    M = \coprod_{w\in \mathcal{R}} B_nwB_m.
  \end{align*}
  \end{theorem}
  
  
  \begin{proof}
    Firstly, we prove that the disjoint union is well-defined i.e. for $w_1,w_2\in \mathcal{R}$, $w_1\neq w_2\Rightarrow B_nw_1B_m \cap B_nw_2B_m = \emptyset$. 
    Suppose $w_1\neq w_2$, if there exist $A_1\in B_nw_1B_m$ and $A_2\in B_nw_2B_m$ such that $A_1=A_2$.
    By Lemma\ref{pr}, we have $A_1 = A_2 \Rightarrow r_{ij}(A_1)=r_{ij}(A_2) \Rightarrow r_{ij}(w_1)= r_{ij}(w_2)$.
    By Lemma\ref{rpr}, every rook matrix can be reconstructed by the partial rank function, then $r_{ij}(w_1)= r_{ij}(w_2)\Rightarrow w_1 = w_2$
    ,which is oppose to the assumption.
    Therefore, the disjoint union is well-defined. 

    Next, we need to proof the existence of decomposition i.e. for any $g\in M$ there exist $w\in \mathcal{R}$ such that
    \begin{align*}
      B_ngB_m = B_nwB_m.
    \end{align*}

    Consider the coset $gB_m$ we can find a representative $g'$ such that 
    \begin{enumerate}[a)]
      \item Each pivot is 1. (pivot is the bottommost nonzero entry of each nonzero column) 
      \item the entries to right of each pivot within its row are all 0, and
      \item there is at most one pivot in each row and at most one pivot in each column. 
    \end{enumerate}
    The algorithm which produces $g'$ from $g$ is: starting with the leftmost nonzero column, find its bottommost nonzero entry, and scale the column
    to make this entry a 1, creating the pivot in this column.Now use this pivot to clear out all entries in its
    row to its right, using column operations that subtract multiples of this column from later columns. Having
    done this, move on to the next nonzero column to the right, and repeat, scaling to create a pivot, and using
    it to eliminate entries to its right.

    A similar algorithm using left multiplication by $B_n$ shows that $f\in B_ngB_m$.

    \qed



  \end{proof}



\newpage

\section{The Białynicki-Birula decomposition of non-square matrices}
In this section we want to apply the Białynicki-Birula decomposition to $M/B_n$. 
Recalled that Białynicki-Birula decomposition is 
\begin{theorem}
  If $X$ be a smooth algebraic variety over $k$ equipped with a locally affine action of $\mathbb{G}_m$.
  Then
    \item $Z$ be a connected component of $X^{\mathbb{G}_m}$. There exist a unique smooth subvariety $X(Z)$ of $X$ such that
    \begin{align*}
      X(Z) = \{y\in X|\lim_{t\to 0} ty\text{ exists and lies in }Z  \}
    \end{align*}
    and unique regular map $\gamma_Z:X(Z)\to Z$ sending $y\in X(Z)$ to the limit $\lim_{t\to 0}ty\in Z$
\end{theorem}

Another important lemma we used is 
\begin{lemma}
  Let $T$ be a torus and $(V,r)$ a finite-dimensional representation of $T$. There exists a cocharacter $\lambda:\mathbb{G}_m\to T$ such that $\mathbb{P}(V)^{\lambda(\mathbb{G}_m)}=\mathbb{P}(V)^T$
\end{lemma}

There are several question we need to answer
\begin{enumerate}
  \item Is $M/B_m$ a projective variety?
  \item Is $M/B_m$ smooth?
  \item Is there exist suitable torus action on $M/B_m$ which is locally affine?
\end{enumerate}
The answer of three question is yes.
We prove it as follow.

Let $k$ be a algebraically closed field, $M$ be the set of $n\times m$ matrices over $k$, 
$B_m$ denotes the invertible upper triangular matrices over $k$ which is the ordinary Borel subgroup of $\gl_n(k)$.
\subsection{$M/B_m$ is a projective variety}

Similar to the $G/B$ which is identified as a flag, we identify $M/B_m$ as a quasi-flag.
First of all, we give the construction of quasi-flag whcih is the generalization of ordinary flag.
\begin{definition}
  Let $V$ be the $n$ dimensional $k$-linear space. Define the $m$ quasi-flag of $V$ by 
  \begin{align*}
    \mathcal{F}(m,V) := \left\{  U_1\subseteq_{\text{linear}} \cdots \subseteq_{\text{linear}} U_m|\text{$U_i\subseteq_{\text{linear}} V$ and $\dim(U_i/U_{i-1})\le 1$ for $i\in\{2,\cdots,m\}$} \right\}.
  \end{align*}
\end{definition}

Next, we give the reason why we can use quasi-flag to identify $M/B_m$.
\begin{lemma}
  $m$-quasi flag is one-to-one corresponding to the $M/B_m$. Hence, the correspondence is given by
  \begin{align*}
    gB_m \mapsto \{U_1\subseteq \cdots \subseteq U_m\}
  \end{align*}
  where $g\in M$ and $U_i =\text{span}\left(\text{ent}_{*1}(g),\cdots,\text{ent}_{*i}(g)\right)$.
\end{lemma}
\begin{proof}
  First, we need to prove it is well-defined i.e. not depends on the choice of representation element $g\in M$.
  Suppose there exists $g'\neq g$ such that $g'B_m=gB_m$, then exists $b\in B_m$ such that $g' = gb$. 
  By definition, the image of $g'B$ is given by
  \begin{align*}
    U'_i = \text{span}\left(b_{11}\text{ent}_{*1}(g),\cdots,b_{1i}\text{ent}_{*1}(g)+\cdots+ b_{1i}\text{ent}_{*i}(g)\right)
  \end{align*}
  where $b_{jj}\neq 0$ for $j\in \mathbb{N}^*_{\le i}$. Since $\det(b)\neq 0$, then
  it is linearly equivalence to the $\{\text{ent}_{*1}(g),\cdots,\text{ent}_{*i}(g)\}$, which implies for $i\in \mathbb{N}^*_{\le m}$, we have 
  $U_{i}=U'_i$. Therefore, it is well-defined.

  Next, we prove the surjective. Pick $u_1\in U_1\backslash \{0\}$ if $\dim(U_1)\neq 0$, otherwise, let $u_1=0$.
  Hence, pick $u_i\in U_{i}/U_{i-1}$ if $\dim(U_{i}/U_{i-1})\neq 0$ or $u_i=0$. Clearly, we have $U_i=\text{span}(u_1,\cdots,u_i)$. So just let 
  $g = [u_1,\cdots,u_m]$, the image of $gB_m$ is the given quasi-flag.

  Finally, the injective is also be satisfied. For fixed quasi-flag $\{U_1\subseteq \cdots \subseteq U_m\}$, we recover the matrix $g$ through the method
  we used in the proof of surjective. For any $g'B$ such that the image is same as the $gB$. Then for any $i\in \mathbb{N}^*_{\le m}$,
  \begin{align}
    \text{span}\left(\text{ent}_{*1}(g),\cdots,\text{ent}_{*i}(g)\right) = \text{span}\left(\text{ent}_{*1}(g'),\cdots,\text{ent}_{*i}(g')\right). \tag{$*$}
  \end{align} 
  ($*$) implies the fact that exists the $\{\lambda_{ji}\}_{j=1}^{i}$ such that 
  \begin{align*}
    \text{ent}_{*i}(g') = \lambda_{1i} \text{ent}_{*1}(g')+\cdots+\lambda_{ii} \text{ent}_{*i}(g')
  \end{align*}
  Suppose that $\lambda_{ii}=0$, then $U_i=U_{i-1}$. Hence, $\text{ent}_{*i}(g)$ can be linear represented by $\{\text{ent}_{*1}(g),\cdots,\text{ent}_{*(i-1)}(g)\}$ i.e.
  exists $\{\mu_j\}_{j=1}^{i-1}$ such that
  \begin{align*}
    \text{ent}_{*i}(g) = \mu_1\text{ent}_{*1}(g)+\cdots \mu_{i-1}\text{ent}_{*(i-1)}(g)
  \end{align*}
  Then 
  \begin{align*}
    \text{ent}_{*i}(g') = (\lambda_{1i}-\mu_1) \text{ent}_{*1}(g')+\cdots+ \text{ent}_{*i}(g')
  \end{align*}
  and let $\lambda'_{ji} = \lambda_{ji}-\mu_j$. If $\lambda_{ii}\neq 0$ then just let $\lambda'_{ji}=\lambda_{ji}$.
  Let $b\in B_m$ with the $ij$ entry $\lambda'_{ij}$ if $i\le j$, then $g'=gb\Rightarrow g'B_m=gB_m$.

  \rightline{$\qedsymbol$}
\end{proof}

Before give the proof of quasi-flag is variety, we can consider another concepts which is more precise than quasi-flag.
\begin{definition}
  Define the Grassmannian $d$-sum of $V$ by
  \begin{align*}
    \tilde{G}_{d,V} = \{U\subseteq_{\text{linear}} V|\dim(U)\le d\}
  \end{align*} 
  the set of linear subspace of $V$ such that the dimension is lower than $d$.
\end{definition}

Obviously, we can embed quasi-flag $\mathcal{F}(m,V)$ in the cartesian product of Grassmannian sum by
\begin{align*}
  \mathcal{F}(m,V) \to \tilde{G}_{1,V} \times \cdots\times \tilde{G}_{\min\{n,i\},V}  \times \cdots \times \tilde{G}_{\min\{n,m\},V}
\end{align*}
Notice that quasi-flag require $\dim(U_{i}/U_{i-1})\le 1$,so $\dim(U_i)\le \min\{i,n\}$
then copies of Grassmannian sum is large enough.

\begin{definition}
  Define the sum Plücker embedding $\tilde{\mathcal{P}}_d: \tilde{G}_{d,V}\to \mathbb{P}(\underline{\bigwedge}^d(V))$ such that for $0\le i\le d$ diagrams 


\newcommand{\tpd}{\tilde{\mathcal{P}}_d}
\newcommand{\pd}{{\mathcal{P}}_d}

\begin{center}
  \begin{codi}[tetragonal]
  \obj[tetragonal=base 6em height 3em]
  { |(A)|G_{{i},V}  & |(B)|  {\bigwedge}^{i}(V)& |(E)| \mathbb{P}({{\bigwedge}}^{i}(V))\\
  &&
  \\ |(C)|\tilde{G}_{{d},V} & |(D)| \mathbb{P}({\underline{\bigwedge}}^{d}(V))& \\ };
  \mor A \rho:-> B \pi:-> D;
  \mor A \iota:-> C;
  \mor C \tpd:-> D;
  \mor B \pi':-> E;
  \mor E \iota:-> D;
  \end{codi}
\end{center}
commute. Where 
\begin{align*}
  \rho:\tilde{G}_{{n},V} &\to {\underline{\bigwedge^{d}}}(V)\\
  W &\mapsto w_1\wedge \cdots \wedge w_{\dim(W)}
\end{align*}
$\{w_1,\cdots,w_{\dim(W)}\}$ is a basis of linear subspace $W$ and $\pi,\pi'$ is the natural homomorphism from affine space to projective space. 
$\iota$ denotes canonical injection.
\end{definition}

\begin{proposition}
  $\tilde{\mathcal{P}}_d|_{G_{{i},V}}(G_{{i},V})$ is isomorphic to the image of Plücker embedding $\mathcal{P}_i(G_{{i},V})$.\label{1}
\end{proposition}
\begin{proof}
  By definition, we can decomposite the Plücker embedding $\mathcal{P}_i =  \pi'\circ\rho $ and $\tilde{\mathcal{P}}_d|_{G_{{i},V}} =  \iota\circ \pi'\circ\rho$
  Since canonical injection does't change the structure of variety i.e. It is regular and biregular because regular function has only $0$ and $\iden$.
  So isomorphic  is been guarantee directly.

  \qed
\end{proof}


\begin{theorem}
  Sum Plücker embedding is a closed embedding in projective space.Therefore, Grassmannian $d$-sum is a projective variety.
\end{theorem}
\begin{proof}
  We first claim that it is well-defined.
  For any linear subspace $U$ of $V$, suppose that $\dim(U) = i$. Then we have
  \begin{align*}
    \tilde{\mathcal{P}}_d(U) =\mathcal{P}_i(U).
  \end{align*}
  Since $\mathcal{P}_i$ is well-defined, $\tilde{\mathcal{P}}_d$ is well-defined i.e. it does not depend on the choice of basis for any linear subspace of $V$.
  Next, If there are linear subspaces $W_1,W_2$ such that $\tilde{\mathcal{P}}_d(W_1)=\tilde{\mathcal{P}}_d(W_2)= \overline{w}$ where  
  \begin{align*}
    \overline{w} = \overline{ w_1\wedge\cdots \wedge w_{l}}
  \end{align*}
  then $\dim(W_1)=\dim(W_2)= l$. According to Proposition \ref{1}, We have
  \begin{align*}
    \tilde{\mathcal{P}}_d(W_i) =\mathcal{P}_l(W_i).
  \end{align*}
  Hence, according to Proposition \ref{2}, $\mathcal{P}_l$ is injective. Then $W_1=W_2$ implies $\tilde{\mathcal{P}}$ is injective.
  Plücker embedding $\mathcal{P}_l $ gives a closed subvariety structure in ${G}_{l,V}$, by \ref{1} we know $\tilde{\mathcal{P}}_d$ give the same structure i.e. there are isomorphic as projective variety.
  By construction of sum Plücker embedding, we see $\tilde{\mathcal{P}}$ gives a closed subvariety structure as well.

  \qed
\end{proof}


\begin{remark}
  Actually, we can give the construction of polynomials on the ${\underline{\bigwedge}^{d}}(V)$ that common locus is the image of sum Plücker embedding.
  \begin{align*}
    \tilde{\mathcal{P}}_d(\tilde{G}_{d,V}) &=\\ &\mathcal{V}\left\{ \text{the $n-i+1$ minors of matrix $\tilde{D}_i,1 \le i \le d$}  \right\}
    \cap \mathcal{V}\left\{ \text{$a_Ia_J=0 |$ for all $|I|\neq |J|$}  \right\}
  \end{align*}
  where $a_I$ denotes Plücker coordinates and $\tilde{D}_i$ is the matrix representation of $\tilde{\varphi}_i$.
\end{remark}

Let we back to the quasi-flag.

\begin{definition}
  Define quasi-flag embedding of $\mathcal{F}(m,V)$ by
  \begin{align*}
    \mathscr{P}: \mathcal{F}(m,V)  &\to \mathscr{S}_m \left(\prod_{i=1}^{m} \mathbb{P} \left(
    \underline{\bigwedge}^{\min(i,n)}(V) \right)\right)   \\
    \{U_1\subseteq \cdots \subseteq U_m\}  &\mapsto \mathscr{S}_m(\overline{u_{11}},\cdots,\overline{u_{m1}\wedge \cdots u_{m\dim(U_m)}} )
   \end{align*}
  where $\{u_{i1},\cdots,u_{i\dim(U_i)}\}$ is a basis of linear subspace $U_i$.
  $\mathscr{S}_m$ denotes Segre embedding of $m$ copies of projective space.
\end{definition}

\begin{theorem}
  Quasi-flag embedding is injective and has closed image in projective space, so quasi-flag is a projective variety. \label{qfp}
\end{theorem}


\begin{proof}
  Proposition\ref{252} implies that if there are two different basis $\{u_i\}$ and $\{{u}'_i\}$ of the $l$-dimensional linear subspace $U$,by the Plücker embedding, $\overline{\{u_1\wedge\cdots\wedge u_l\}} = \overline{\{u'_1\wedge\cdots\wedge u'_l\}}$.
  So the embedding of quasi-flag is well-defined as well.
  Secondly, we need to porve that it is injective. Let 
  \begin{align*}
    \mathscr{P}':\mathcal{F}(m,V)&\to  \prod_{i=1}^{m} \mathbb{P}(\underline{\bigwedge}^{\min(i,n)}(V)) \\
    \{U_1\subseteq \cdots \subseteq U_m\}  &\mapsto (\overline{u_{11}},\cdots,\overline{u_{m1}\wedge \cdots u_{m\dim(U_m)}} )
  \end{align*}
  If exists $U,U'\in \mathcal{F}(m,V)$ such that $\mathscr{P}'(U)=\mathscr{P}'(U')=\alpha$, then for all $i\in \{1,\cdots,m\}$ 
  \begin{align*}
    U_i=U'_i = \mathcal{P}_{\min\{i,n\}}^{-1}\circ \pi_{i} (\alpha) \ \tag{$*$}
  \end{align*}
  where $\pi_i$ denote the ordinary projection of $m$-tuples. ($*$) shows that $\mathscr{P}'$ is injective.
  Recall that Segre embedding is also injective, then $\mathscr{P} = \mathscr{S}_m\circ \mathscr{P}'$ is injective.

 

  Following the definition, there are only two conditions that quasi-flag need.
  \begin{enumerate}[a)]
    \item $U_{i-1}\subseteq U_{i}$
    \item $\dim(U_i/U_{i-1}) \le 1$ for $i=\{2,\cdots,m\}$
  \end{enumerate}
  It is clear that quasi-flag is embeded in the $m$-tuples of Grassmannian sum which is closed because Segre embedding is morphism of projective varieties. 
  All we need to prove is that quasi-flag is a closed subvariety of $m$-tuples of Grassmannian sum.
  From the discussion of exterior power, 
  \begin{align*}
    U_{i-1}\subseteq U_{i}\Leftrightarrow u_{i1}\wedge\cdots u_{i\dim(U_i)}\wedge u_{(i-1)j} \tag{$*$}
  \end{align*} 
  where $\{u_{ij}\}_{j=1}^{\dim(U_i)}$ is a basis of $U_i$ and $\{u_{(i-1)j}\}_{j=1}^{\dim(U_{i-1})}$ is a basis of $U_{i-1}$. 
  Hence, because it is totally decomposable, then we consider ($*$) is also equivalence to 
  \begin{align*}
    B_{u_i} u_{(i-1)j}=0 ,j\in\{1,\cdots,\dim(U_{i-1})\}
  \end{align*}
  Which are homogeneous polynomials.

  For condition $b)$ Let $a_{Ij}$ denotes the Plücker coordinate $a_I$ in $j$th projective space. Then, the homogeneous coordinate is given by
  \begin{align*}
  [a_{\emptyset 1}\times\cdots\times a_{\emptyset m}:\cdots:a_{{I_1}1}\times\cdots\times a_{{I_m}m} :\cdots]
  \end{align*}
  If one entry $a_{{I_1}1}\times\cdots\times a_{{I_m}m} $
  of homogeneous coordinate not equals to 0, then it can determine the rank sequence $ \{\rank{U_i}\}_{i=1}^n$
  by $\rank{U_i} = |I_i|$, then we let the polynomials 
  \begin{align*}
    a_{{I_1}1}\times\cdots\times a_{{I_m}m} = 0
  \end{align*}
  If exists $i\in \mathbb{N}^*_{\le m-1}$ such that $|I_{i+1}|-|I_{i}|>1$.
  Therefore, quasi-flag is a projective varieties.
\end{proof}

\subsection{The smoothness of $M/B_m$}
Suppose the quasi-flag embedding  send $M/B$ to the projective space $\mathbb{P}(W)$.
Our strategy of proving smoothness is find the appropriate smooth open cover of $M/B$ on $\mathbb{P}(W)$.
Before find the cover, we give the proof of some useful lemma.

\begin{definition}
  Define the extended flag  of $V$ as:
  \begin{align*}
    \mathcal{F}(n_1,\cdots,n_r;V) = \{U_1\subseteq \cdots\subseteq U_r|U_i\subseteq_{\text{linear}} V ,\dim(U_i)=n_i\}
  \end{align*}
  where $1\le n_1\le\cdots \le n_r$ is a sequence of integer.
\end{definition}

\begin{lemma}
  There exist a flag variety $\mathcal{F}(j_1,\cdots,j_q;V)$ such that it is isomorphic to 
  extended flag $ \mathcal{F}(n_1,\cdots,n_r;V)$. \label{exfs}
\end{lemma}
\begin{proof}
  Obviously holds by construction. 
  \qed
\end{proof}

\begin{theorem}
  $M/B_m$ is smooth.
\end{theorem}
\begin{proof}
  Recall that the proof of Theorem\ref{qfp}, we represent quasi-flag which is identified to $M/B_m$ as the locus set of a group of polynomials. 
  under the closed subvariety $\mathscr{S}(\tilde{G}_{1,V}\times\cdots\times \tilde{G}_{\min(n,m),V}) $. The construction of polynomials as follow.
  \begin{align*}
    &\{B_{u_i} u_{(i-1)j}=0 ,j\in\{1,\cdots,\dim(U_{i-1})\} \} \\
    \cup&\{a_{{I_1}1}\times\cdots\times a_{{I_m}m} = 0|\text{If exists $i\in \mathbb{N}^*_{\le m-1}$ such that $|I_{i+1}|-|I_{i}|>1$}\}
  \end{align*}
  We consider the closed subvariety 
  \begin{align*}
    \mathscr{S}(\tilde{G}_{1,V}\times\cdots\times \tilde{G}_{\min(n,m),V})  \cap 
    \mathcal{V}(B_{u_i} u_{(i-1)j}=0 ,j\in\{1,\cdots,\dim(U_{i-1})\})
  \end{align*}
  which includes the quasi-flag, and consider part of affine chart  
  \begin{align*}
    \{a_{{I_1}1}\times\cdots\times a_{{I_m}m} \neq 0|\text{$|I_i|$ is increasing} \}.
  \end{align*}
  It is a open cover of quasi-flag in a closed subvariety. Hence each open cover is corresponding to 
  a extended flag variety because it satisfied the including condition i.e. 
  \begin{align*}
    \mathcal{V}(B_{u_i} u_{(i-1)j}=0 ,j\in\{1,\cdots,\dim(U_{i-1})\})
  \end{align*}
  Then we can recover the extended flag $\mathcal{F}(|I_1|,\cdots,|I_m|;V)$ which is smooth according to Lemma\ref{exfs}.
\end{proof}


\subsection{The torus action on $M/B_m$}

Recall that the ordinary torus action on algebraic group is $t^{-1}gt$ where $t\in T$ and $g\in G$.
However, There are a few question when we consider the torus action $M/B_m$. The most important thing is that we can't find the trous directly from $M$.
All we can do is just keep the consistence in form.
 
We define the torus action and $\mathbb{G}_m$ of $T$ firstly.
\begin{definition}
  Suppose $M$ is the $n\times m$ matrix over filed $k$, $T_j$ be the $j$-dimensional invertible diagonal matrix over $k$.
  Define torus $T_{n+ m}$ action on $M$ by 
  \begin{align*}
    T_{n+ m} \times M &\to M\\
    t\times g &\mapsto t_ngt_m
  \end{align*}
  where $t=t_n\times t_m\in T_{n+m}$ and $g\in M$.
  
  Define co-character $\lambda: \mathbb{G}_m \to T_n\times T_m$ such that 
\begin{align*}
  h \mapsto \text{diag}(h^{x_1},\cdots,h^{x_n}) \times \text{diag}(h^{y_1},\cdots,h^{y_m})
\end{align*}

where $x_i,y_i\in \mathbb{Z}$. Define the $\mathbb{G}_m$-action on $M$ through co-character $\lambda$ by
\begin{align*}
  h.g = \lambda_1(h)g\lambda_2(h)
\end{align*}
where $\lambda_i = \pi_i\circ\lambda$ and $\pi_i$ denotes projection.
\end{definition}

Because $M$ is a affine variety then the torus action is clearly locally affine.
Further more, we want to give the torus action on the $M/B_m$.

\begin{definition}
  Define $T_n\times T_m$ torus action on $M/B_m$ by
  \begin{align*}
    t \times gB_m &\mapsto (t_ngt_m)B_m
  \end{align*}
  where $t=t_n\times t_m\in T_{n+m}$ and $g\in M$.
  Define the $\mathbb{G}_m$-action on $M/B_m$ through co-character $\lambda$ by
  \begin{align*}
    h.gB_m = \lambda_1(h)g\lambda_2(h) B_m
  \end{align*}
  where $\lambda_i = \pi_i\circ\lambda$ and $\pi_i$ denotes projection.
\end{definition}


\begin{theorem}
  $\mathbb{G}_m$-action is locally affine on $M/B_m$.
\end{theorem}
\begin{proof}
  Consider the quasi-flag embedding of $M/B_m$, then the projective space we choice have the homogeneous
  coordinates 
  \begin{align*}
    [a_{\emptyset 1}\times\cdots\times a_{\emptyset m}:\cdots:a_{{I_1}1}\times\cdots\times a_{{I_m}m} :\cdots]
  \end{align*} 
  with the affine chart $\{U_I|I=I_1\times \cdots \times I_m,I_i\subseteq \mathbb{N}_{\le i}^* \}$ which is isomorphic to affine space.
  We choose part of affine chart whcih satisfy the increasing condition $|I_{i+1}|-|I_{i}|\le 1$. It is a open cover of $M/B_m$.
  Finally,we claim that it is a $\mathbb{G}_m$-invarient. For fixd affine chart $U_I$ we have
  \begin{align*}
    a_{{I_1}1}\cdots a_{{I_m}m} \neq 0 \tag{$*$}
  \end{align*}

  Suppose $I\not=\{\emptyset\}\times \cdots \times\{\emptyset\}$, let $j_0=\min_{j}\{I_j\neq \emptyset\}$, $J^* = \{i||I_{i}|-|I_{i-1}|=1,i\ge j_0+1\} =\{j_1<\cdots<j_q\}$ and $J=\{j_0<\cdots<j_q\}$.
  then for any $gB_m\in\mathscr{P}^{-1}(U_I)$ 
  satisfied the conditions
  \begin{enumerate}[a)]
    \item $\text{ent}_{*l}(g) = 0 $ for $1\le l\le j_0 $
    \item $|\Delta_{I_i,J(i)}|\neq 0$ where $i\ge j_0,J(i) = \{j_0,\cdots,j_{|I_i|-1}\}$
    \item For $i\notin J$, 
    let $J'(i) = \{j_0,\cdots,j_l,i\}$ where $l=\max_{l}\{j_l\le i\}$.
    Any maximal minors $\Delta$ of $\Delta_{\mathbb{N}^*_{\le n},J'(i)}$ such that $|\Delta|=0$.
  \end{enumerate}

  Hence, it is sufficient. Suppose $gB_m\in M/B_m$ satisfy those conditions.
  And consider the quasi-flag embedding in product of projective spaces
  \begin{align*}
    \mathscr{P}':\mathcal{F}(m,V)&\to  \prod_{i=1}^{m} \mathbb{P}(\underline{\bigwedge}^{\min(i,n)}(V))
  \end{align*}
  and $\tilde{g} = \mathscr{P}'(gB_n)$, $\pi_j:\prod_{i=1}^{m}\mathbb{P}(\underline{\bigwedge}^{\min(i,n)}(V))\to \mathbb{P}(\underline{\bigwedge}^{\min(j,n)}(V))$ is a projection.
  Then a) implies $\pi_{l}(\tilde{g}) = a_{\emptyset l}$. Condition b) give a linear independent set of vector 
  \begin{align*}
    A = \{\text{ent}_{*j_0},\cdots,\text{ent}_{*j_q}\} 
  \end{align*}
Hence it is a basis of $g$ according conditions c) which implies that $A$ is maximally linearly independent.
Then we have  $\pi_i(\tilde{g})\in \mathbb{P}({\bigwedge}^{l}(V))$, where $l=|\{j\le i|j\in J\}|$. For any $j_l \le i< j_{l+1}$ and $|I_k|=l$, 
The Plücker coordinate in product space $a_{I_k,i}\neq 0$ which implies $a_{{I_1}1}\cdots a_{{I_m}m} \neq 0$.
Then $g$ satisfy condition a),b),c) if and only if $gB_m\in\mathscr{P}^{-1}(U_I)$.

It is easy to see that all the conditions is invarient of torus action because it not change $|\Delta|$ equals to $0$ or not.

Finally, $\mathscr{P}^{-1}(U_I)$ is affine by construction of polynomials which is the determinnant minors of matrix.

\qed
\end{proof}





\newpage 
\bibliographystyle{siam}
\bibliography{references}

\end{document}